%% ===================================================================
%% Missense pathogenicity prediction in the DNA mismatch repair genes
%% Target: Oxford Bioinformatics (Original Paper)
%%
%% BUILD
%%   latexmk -pdf main.tex        (or: pdflatex; bibtex; pdflatex x2)
%%
%% bioinfo.cls is distributed by OUP and is NOT in this repository. If it is
%% absent the preamble falls back to `article` with matching no-op macros, so
%% the document still builds for local review; the typeset result will not be
%% the journal's layout. See README.md in this directory.
%%
%% EVERY NUMBER IN THIS FILE IS TRACEABLE TO A REPOSITORY ARTIFACT.
%% Unverified items appear as \todo{...} and are enumerated in
%% ../../MISSING_EVIDENCE.md. Do not replace a \todo with a plausible value.
%% ===================================================================

%% The class test only sets a flag. The fallback definitions must NOT sit
%% inside an \IfFileExists branch: that branch is a macro argument, so a `#1'
%% written there is read as a parameter of the argument and every definition
%% taking one errors with "Illegal parameter number in definition of
%% \reserved@a". LaTeX recovers by assuming `##', which is why this built at
%% all, but three errors per run is not a state to leave a submission in.
\IfFileExists{bioinfo.cls}{\def\bi@hasclass{1}}{\def\bi@hasclass{0}}

\if1\bi@hasclass
  \documentclass{bioinfo}
  \copyrightyear{2026}
  \pubyear{2026}
  \access{Advance Access Publication Date: Day Month Year}
  \appnotes{Original Paper}
\else
  % ---- Fallback: no bioinfo.cls available -------------------------------
  \documentclass[11pt,a4paper]{article}
  \usepackage[margin=2.4cm]{geometry}
\fi

\if1\bi@hasclass\else
  \makeatletter
  %% article's \title and \author take no optional argument. Left unwrapped,
  %% TeX grabs the single token `[' as the whole argument: the running head and
  %% the real title are then set as body text above the title block, and
  %% \maketitle prints a lone `[' for each. Swallow the optional argument.
  \let\bi@title\title
  \renewcommand{\title}[2][]{\bi@title{#2}}
  \let\bi@author\author
  \renewcommand{\author}[2][]{%
    \bi@author{\parbox{0.92\textwidth}{\centering #2}}}
  %% bioinfo.cls typesets these inside the title block. article has nowhere to
  %% put them, and printing on the spot puts them above the title and leaves
  %% \maketitle to emit an empty block on the next page. Collect instead, and
  %% let \maketitle place them.
  \newcommand{\bi@front}{}
  \newcommand{\bi@abstract}{}
  \newcommand{\copyrightyear}[1]{}
  \newcommand{\pubyear}[1]{}
  \newcommand{\access}[1]{}
  \newcommand{\appnotes}[1]{}
  \newcommand{\firstpage}[1]{}
  \newcommand{\subtitle}[1]{}
  \newcommand{\history}[1]{}
  \newcommand{\editor}[1]{}
  \newcommand{\address}[1]{\g@addto@macro\bi@front{\par #1}}
  \newcommand{\corresp}[1]{\g@addto@macro\bi@front{\par #1}}
  \renewcommand{\abstract}[1]{\gdef\bi@abstract{#1}}
  \date{}
  \let\bi@maketitle\maketitle
  \renewcommand{\maketitle}{%
    \bi@maketitle
    {\centering\small\bi@front\par}%
    \begin{quote}\small\bi@abstract\end{quote}}
  \makeatother
\fi

\usepackage{booktabs}
\usepackage{array}
\usepackage{tabularx}
%% Weighted, ragged-right X column. The weights within one table must sum to
%% the number of X columns it uses. Ragged rather than justified: at these
%% widths justification opens holes between words.
\newcolumntype{L}[1]{>{\hsize=#1\hsize\raggedright\arraybackslash}X}
%% Font expansion and margin protrusion. Guarded so a TeX installation without
%% it still builds.
\IfFileExists{microtype.sty}{\usepackage{microtype}}{}
\usepackage{graphicx}
%% Figures live where scripts/make_figures.py writes them. Keeping a second
%% copy in this directory only invites the two to drift, so the submission
%% bundle should be assembled by copying ../figures at the time it is built.
\graphicspath{{./}{../figures/}}
\usepackage{amsmath}
\usepackage{siunitx}
\usepackage[table]{xcolor}
\usepackage[colorlinks=true,allcolors=blue]{hyperref}

\sisetup{detect-weight=true, group-minimum-digits=4}

%% Let long URLs break rather than overrun the measure.
\Urlmuskip=0mu plus 1mu\relax

%% Unresolved evidence. Deliberately loud: a placeholder that reads like prose
%% is how a fabricated number reaches a reviewer.
\newcommand{\todo}[1]{\textcolor{red}{\textbf{[TODO: #1]}}}

%% Gene symbols are italicised; protein products are not.
\newcommand{\gene}[1]{\textit{#1}}

\firstpage{1}

\begin{document}

\title[Held-out-gene evaluation of missense pathogenicity models]{%
Merging heterogeneous variant evidence for missense pathogenicity prediction:
a leakage-controlled benchmark in the DNA mismatch repair genes}

\author[\todo{Author} \textit{et~al}.]{%
\todo{Author list --- do not invent; complete from the actual contributors}}

\address{\todo{Affiliations}}

\corresp{\todo{Corresponding author contact}}

\history{Received on XXXXX; revised on XXXXX; accepted on XXXXX}

\editor{Associate Editor: XXXXXXX}

\abstract{
\textbf{Motivation:}
Evidence bearing on a missense variant is scattered across clinical archives,
deep mutational scanning (DMS) assays, population frequency databases,
structural models and precomputed pathogenicity scores. These sources disagree,
overlap, and in several cases encode the label a classifier is asked to
predict: population allele frequency is itself a benignity criterion, and
precomputed scores were trained on the same clinical archives used for
evaluation. Reported accuracy is further inflated by random splits, under which
different substitutions at one residue appear on both sides of the split. What
is missing is not another model but a merge and evaluation protocol under which
a reported number means something.
\\
\textbf{Results:}
We assembled a variant-level evidence table for \gene{MLH1}, \gene{MSH2},
\gene{MSH6} and \gene{PMS2}, keyed on
\texttt{(uniprot\_id, position, wt\_aa, mut\_aa)} with the canonical UniProt
sequence as the sole coordinate authority. The row universe is the union of all
possible missense substitutions (\num{74328} variants; \num{3912} residues
$\times$ 19), not an anchor on any one source. Labels are resolved only after
merging, under an explicit precedence, and contradictory evidence is
quarantined rather than coerced. Evaluating by leave-one-gene-out, we find that
full fine-tuning of a 650M-parameter protein language model does not improve
ranking over a frozen backbone once the feature set is held constant
($\Delta$AUROC $-0.007$, $\Delta$AUPRC $-0.005$; three seeds), whereas adding 27
curated prior features raises AUROC by $0.037$. A gap previously attributed to
freeze depth is attributable to the feature set. Fine-tuning nonetheless shifts
the operating point substantially (accuracy $0.698 \rightarrow 0.840$, recall
$0.702 \rightarrow 0.872$), which AUROC alone does not reveal. A curated-feature
head reading no language-model input at all scored $0.9507$ mean AUROC over the
three scoreable genes --- above every cell in the grid, including those reading
the same 27 features. Feature-family ablations at three seeds localise that
contribution: removing the 18 external prior-score columns costs $0.038$ AUROC
and removing them together with allele frequency costs $0.059$, while removing
structural or domain features moves the mean by $+0.004$, inside seed noise.
Threshold selection on an inner validation fold beat a fixed $0.5$ cut by only
$0.007$ accuracy and was worse in 6 of 16 configurations, indicating that
threshold transfer across held-out genes, not model capacity, is the binding
constraint.
\\
\textbf{Availability and implementation:}
Code, build manifests with per-source checksums, and all per-variant
predictions are at
\url{https://github.com/jvikramsrd/variant_pathogenicity_dl}.
\todo{public release tag}
\\
\textbf{Contact:} \todo{email}
\\
\textbf{Supplementary information:} Supplementary data are available at
\textit{Bioinformatics} online.}

\maketitle

% =====================================================================
\section{Introduction}
% =====================================================================

Clinical sequencing returns more missense variants than can be interpreted. In
hereditary cancer panels the dominant output is the variant of uncertain
significance (VUS), which by definition cannot guide management. Computational
predictors scale where curation does not, and several are now embedded in
variant-curation workflows. The difficulty is establishing what a reported
accuracy figure means.

Three problems recur. First, the evidence for a variant is distributed across
sources with incompatible semantics. A clinical archive records an asserted
significance with a review status; a DMS experiment records a continuous
fitness measurement under one assay in one cell type; a population database
records how often an allele is observed. Combining them requires deciding which
are evidence, which are features, and which are labels --- a decision that is
usually implicit and therefore unauditable.

Second, several sources are label-adjacent. Population allele frequency is a
benignity criterion in clinical classification guidelines~\cite{richards2015acmg};
if allele frequency both mints benign labels and enters the feature matrix, a
classifier recovers the label by construction. Precomputed scores such as
AlphaMissense were trained on population and clinical data, so using them as
features while evaluating against clinical labels risks measuring agreement
with an upstream model rather than with biology~\cite{grimm2015circularity}. These are not hypothetical
failure modes; they are the default outcome of naive merging.

Third, the evaluation split determines what is measured. Under a random split,
variants at the same residue of the same protein routinely appear on both
sides. Because a protein language model encodes local sequence context, and
because curated features such as domain membership are constant within a
region, a model can score a held-out variant using information from its
neighbours. The resulting figure describes interpolation within known genes,
not generalisation to a new one.

We address the merge and evaluation problem for the four DNA mismatch repair
(MMR) genes associated with Lynch syndrome. The gene set is deliberately small:
it is well characterised, carries a substantial ClinVar record and a large DMS
dataset for \gene{MSH2}, and includes \gene{PMS2}, whose pseudogene homology
makes short-read calls unreliable over part of the coding sequence and so
forces an explicit exclusion policy rather than a silent one.

Our contributions are (i) a variant-level merge protocol with a single identity
key, canonical-sequence validation of every wild-type residue, union-based row
construction and per-source provenance; (ii) an explicit label-resolution policy
with documented DMS direction mapping and quarantine of contradictory evidence;
(iii) a leakage-controlled evaluation protocol with leave-one-gene-out as the
primary split; (iv) a 16-cell ablation grid that separates backbone freeze depth
from feature set, two factors previously confounded in our own earlier analysis;
and (v) an account of what this cohort cannot support, including a held-out gene
too small to score and a threshold protocol that transfers poorly across genes.

% =====================================================================
\section{Methods}
% =====================================================================

\subsection{Data sources}

Sources, access URLs and per-file SHA-256 checksums are recorded in the build
manifest (\path{data/mmr/processed/extended/manifest.json}). ClinVar~\cite{landrum2018clinvar}
(\texttt{variant\_summary.txt.gz}) supplies clinical labels; only assertions at
two review stars or more were eligible for supervision. ProteinGym v1.3~\cite{notin2023proteingym}
supplies DMS fitness scores with an assay-binarised label, a secondary clinical
benchmark, and precomputed zero-shot model scores. AlphaMissense v1.0~\cite{cheng2023alphamissense}
(2023-09-18; CC BY-NC-SA 4.0) supplies a pathogenicity prior. gnomAD v4~\cite{chen2024gnomad,karczewski2020gnomad} supplies
allele frequencies and gene-level constraint. UniProt~\cite{uniprot} supplies the canonical
sequence --- the coordinate authority for the entire table --- plus domain and
functional-site annotation. AlphaFold DB~\cite{varadi2022afdb,jumper2021alphafold} (CC-BY-4.0) supplies per-residue
pLDDT, and InterPro~\cite{interpro} (CC0) domain and family intervals. Adapters for MaveDB~\cite{esposito2019mavedb} and
CIMRA/OddsPath exist and are exercised by the test suite, but the manifest for
this build records no source block for either, indicating they were not
enabled \todo{confirm and state explicitly in Table~1}. Per-source counts are
given in Table~\ref{tab:sources}.

\begin{table*}[t]
\caption{Data sources for the MMR panel. Counts are read from
\texttt{manifest.json} for the build identified by SHA-256
\texttt{78eb5d60\ldots}; ``retained'' is the number of panel rows the source
contributed evidence to, not the size of the download. The final column
distinguishes sources that supply supervision from those that supply features
only --- a distinction that is usually left implicit and is the reason several
of these can leak. gnomAD allele frequency covers a minority of rows by
construction: the row universe is every possible substitution, whereas gnomAD
records only observed alleles.}
\label{tab:sources}
\centering
\footnotesize
\setlength{\tabcolsep}{4pt}
%% Five `l' columns of prose ran ~30% past the text block, clipping the last
%% two. tabularx holds the table to \textwidth in either class and wraps the
%% two descriptive columns instead.
\begin{tabularx}{\textwidth}{@{}L{0.95} L{0.80} L{1.25} L{1.15} L{0.85}@{}}
\toprule
Source & Role & Retained on panel & Mapping onto the key & Used for \\
\midrule
ClinVar ($\geq$2\,\textsc{star}) & Clinical assertion & 464 labelled; \num{9589} VUS kept unlabelled & HGVS p.\ $\rightarrow$ UniProt, wt validated & Label (precedence 1) \\
ProteinGym DMS v1.3 & Functional assay & \num{16749} rows, 1 assay (\gene{MSH2}) & UniProt coordinates & Label (precedence 3); not used here \\
ProteinGym clinical v1.3 & Clinical benchmark & 564 joined & RefSeq NP $\rightarrow$ exact sequence match & Label (precedence 2) \\
Zero-shot scores v1.3 & Prior score & 564 joined & UniProt coordinates & Feature \\
AlphaMissense v1.0 & Prior score & \num{74328} & UniProt coordinates & Feature \\
gnomAD v4 & Population frequency & \num{6492} with joint AF; constraint on \num{74328} & Protein substitution & Feature \\
UniProt & Coordinate authority & 4 accessions; 20 domain intervals; 61 sites & --- (defines the key) & Feature and key \\
AlphaFold DB & Structure confidence & \num{3912} residues, 4 accessions & Accession $+$ position & Feature \\
InterPro & Domain/family & 123 intervals; \num{72599} rows in a domain & Accession $+$ position & Feature \\
MaveDB & Functional assay & \todo{0 --- confirm} & --- & Not enabled in this build \\
CIMRA / OddsPath & Calibrated evidence & \todo{0 --- confirm} & --- & Not enabled in this build \\
\bottomrule
\end{tabularx}
\end{table*}

\begin{table}[t]
\caption{Cohort and label composition. The supervised cohort is the 683
clinically labelled variants; the \num{16420} DMS-proxy labels all derive from a
single \gene{MSH2} assay and so contribute no cross-gene supervision under
leave-one-gene-out.}
\label{tab:cohort}
\centering
\small
\begin{tabular}{lr}
\toprule
Quantity & Value \\
\midrule
Master rows (\num{3912} residues $\times$ 19) & \num{74328} \\
Labelled (manifest) & \num{17124} \\
\quad benign / pathogenic & \num{14500} / \num{2624} \\
Clinical labels used for supervision & 683 \\
\quad ClinVar / ProteinGym clinical & 443 / 240 \\
\quad \gene{MLH1} / \gene{MSH2} / \gene{MSH6} / \gene{PMS2} & 208 / 335 / 119 / 21 \\
DMS-proxy labels (all \gene{MSH2}, one assay) & \num{16420} \\
ClinVar VUS retained, unlabelled & \num{9589} \\
\gene{PMS2} excluded by homology gate & 21 \\
Variants with gnomAD joint AF & \num{6492} \\
Variants with gnomAD gene constraint & \num{74328} \\
\bottomrule
\end{tabular}
\end{table}

\subsection{Variant normalisation and master-table assembly}

Every variant is identified by
\texttt{(uniprot\_id, position, wt\_aa, mut\_aa)}. Gene symbol is derived, not
authoritative: source-specific aliases were previously able to collide on the
key, so the panel symbol is written onto every row after the join and the
relabelling count recorded. For every incoming row the wild-type residue is
checked against the canonical sequence at that position; mismatches are dropped
and counted, never coerced (this build: zero dropped from both ClinVar and
DMS). ProteinGym clinical records, keyed by RefSeq protein accession, are
resolved by exact whole-protein sequence matching rather than by name; three
panel accessions matched and \num{2522} did not, the latter excluded and
counted.

The master table is a union of validated sources rather than an anchor on one.
Here the union is the complete saturation set --- \num{3912} residues
$\times$ 19 substitutions $=$ \num{74328} rows --- so every possible missense
substitution is present whether or not any source has an opinion, and
unlabelled variants survive to be scored. Metadata are backfilled before
deduplication so identity is complete first, and export is blocked by a
validator that refuses a table containing a duplicate key, a null key
component, a non-binary resolved label, or a resolved label on a quarantined
row. Sources are attached by deterministic left join; missing values are left
missing and represented by explicit indicators where a model consumes them.
Per-assay DMS records are preserved in a long-format table with only a summary
in the master table, which is what prevents multi-assay genes from creating
duplicate rows.

\paragraph{A provenance defect found and corrected.}
The manifest was originally written immediately after the table, while a later
stage joined gnomAD and rewrote the file. The manifest therefore described the
pre-join table, recording \texttt{include\_gnomad: false} and an artefact
checksum for a file that no longer existed. We report this because a checksum
that does not match the file it names is worse than none: it invites a reader
to trust it. The builder now re-stamps the manifest after any post-hoc
modification, and a repair utility corrects earlier builds in place by measuring
the table rather than trusting recorded flags. The manifest accompanying these
results has been repaired; its digest for \texttt{extended\_dataset.csv} now
agrees with both the file and an independent pre-run capture.

\subsection{Label construction and conflict handling}

Labels are resolved only after all sources are merged, so no source can
pre-empt another by join order. Precedence is (1) ClinVar pathogenic/likely
pathogenic versus benign/likely benign at $\geq$2 stars, (2) ProteinGym
clinical, (3) eligible single-assay DMS proxy.

ProteinGym's \texttt{DMS\_score\_bin} marks the \emph{top} half of assay fitness
--- functionally tolerated --- as 1, so pathogenicity supervision requires
inversion. This is the most dangerous line in the label code: an unnoticed sign
error would invert every DMS-derived label while leaving all summary statistics
superficially plausible. It is documented at the point of use and covered by a
direction test. A DMS proxy label is admitted only when exactly one assay covers
the variant.

VUS are never labels; their review metadata is retained for prospective
scoring. Within a source, a key carrying both binary labels is not resolved by
archive order but excluded from that source's supervision. Across sources,
ClinVar/ProteinGym-clinical disagreement sets a conflict flag which nulls the
resolved label, clears the label source and zeroes the label weight, while raw
per-source columns are preserved. Clearing the label source is deliberate:
leaving a source name on a row whose label was withheld invites downstream code
that groups by source, without checking the conflict flag, to count evidence
excluded on purpose. In the broad 79-gene build, ClinVar and ProteinGym clinical
overlap on \num{1258} variants with zero disagreements.

Allele-frequency-derived priors are quarantined from the feature set whenever
allele frequency also mints benign labels, since that combination makes the
ACMG BS1 flag equal the label by construction. Gene-constant priors are dropped
under leave-one-gene-out, because a feature with one value per gene is a gene
identifier. The runs reported here did not activate frequency-derived labels,
so that quarantine was not engaged; the gene-constant drop was active.

\subsection{Models and evaluation}

We evaluate three model families: a head on the 27 curated prior features with
no ESM input; \texttt{esm2\_t33\_650M\_UR50D}~\cite{lin2023esm2} frozen with a trained head; and
the same backbone fine-tuned, either fully or with the last two layers
unfrozen. The backbone is used in a siamese configuration comparing wild-type
and substituted sequences, with a pseudo-log-likelihood-ratio (PLLR) term~\cite{meier2021pllr}
enterable as a residual, concatenated, or disabled. Optimisation used
micro-batch 1 with 8 accumulation steps (effective batch 8), bfloat16 autocast
and gradient checkpointing, which is what allows a 650M full fine-tune to fit a
16\,GiB card. Early stopping used inner-validation AUROC with patience 3,
restoring the best epoch's weights before the held-out gene was scored.

The primary split is leave-one-gene-out. Within the training genes, inner folds
are grouped on \texttt{"\{uniprot\_id\}:\{position\}"} so different
substitutions at one residue cannot straddle train and validation; grouping on
bare position was a defect in an earlier version of this code, which pooled
position 42 of one protein with position 42 of another. The decision threshold
is chosen by maximising MCC on the inner-validation fold; the held-out gene is
never used for selection, calibration or thresholding. Confidence intervals are
percentile bootstrap with \num{10000} resamples, stratified by class. The
positive class is pathogenic.

Metrics are computed by a single module so every arm is scored by identical
code, and comprise AUROC, AUPRC, accuracy, F1 for the pathogenic class, macro
and weighted F1, balanced accuracy, precision, recall, specificity, negative
predictive value, Brier score, expected calibration error and the confusion
matrix. A cohort is scored only if it holds at least 20 variants and at least 5
of each class; cohorts failing this are reported as unavailable with a reason
rather than scored. This rule excludes \gene{PMS2} from every performance table
below.

Each grid cell writes a results CSV, a per-variant predictions CSV, an
inner-validation predictions CSV and a summary JSON, all tagged with a
deterministic cell slug. Every summary carries a provenance block: the SHA-256
of the table actually read, an order-independent hash of the resolved feature
schema, a hash of the held-out key assignment --- taken over the keys rather
than the counts, since two splits of equal size holding different variants are
different splits --- the git commit with a dirty flag, and the versions of
\texttt{torch}, \texttt{transformers}, \texttt{scikit-learn}, \texttt{numpy}
and \texttt{pandas}. Aggregation across cells is gated on the dataset and split
hashes: pooling refuses runs computed on different tables or different splits,
and additionally requires an identical feature schema among seed replicates of
one arm. Feature schema is deliberately excluded from the pooling gate, because
an ablation arm differs from its comparator in precisely that field and a gate
forbidding it would forbid the ablation table. A run carrying no provenance
block is reported as missing rather than treated as matching.

Post-hoc calibration was \emph{not} applied. Temperature scaling~\cite{guo2017calibration}
and isotonic regression~\cite{zadrozny2002isotonic} are implemented, but the fine-tuning script discarded its
inner-validation probabilities after using them to select a threshold, leaving
no leak-free split on which to fit a calibrator; fitting on the held-out gene is
the leak the protocol forbids. This is corrected in the code but postdates the
runs reported here.

% =====================================================================
\section{Results}
% =====================================================================

\subsection{Dataset construction}

The merge produced \num{74328} variants across four genes, equal to the
complete saturation set, so no variant was lost or duplicated. Export
validation passed and wild-type validation dropped zero rows. The manifest
records \num{17124} labelled variants (\num{14500} benign, \num{2624}
pathogenic) before the pseudogene homology gate is applied. The delivered table
holds \num{17103} (\num{14488} benign, \num{2615} pathogenic), which is exactly
what the label sources sum to: \num{16420} DMS-proxy, 443 ClinVar and 240
ProteinGym clinical. The 21-row difference is the homology-gated \gene{PMS2}
rows (12 benign, 9 pathogenic), which the gate removes from the table rather
than flagging in place --- every surviving row carries
\texttt{pms2\_homology\_excluded == 0}. The manifest count is therefore
pre-gate and the table count post-gate; both are correct and they are not the
same quantity.

All \num{16420} DMS-proxy labels come from a single \gene{MSH2} assay. Under
leave-one-gene-out they therefore contribute nothing when \gene{MSH2} is held
out, and single-gene, single-assay signal otherwise. The supervised cohort for
all model results is consequently the 683 clinically labelled variants
(Table~\ref{tab:cohort}): 208
\gene{MLH1}, 335 \gene{MSH2}, 119 \gene{MSH6}, 21 \gene{PMS2}. A further 21
\gene{PMS2} variants were excluded by the pseudogene homology gate over codons
382--862, derived from MANE Select transcript ENST00000265849~\cite{mane} and self-validated
against the 862-residue canonical sequence.

A broad 79-gene build (\num{1152863} rows $\times$ 68 columns) passes all 11
audit checks and holds \num{189006} labelled variants, but \num{185600} of those
are single-assay DMS and only \num{2148} are clinical. No model was trained on
it \todo{broad-panel model results; see MISSING\allowbreak\_EVIDENCE.md item 4}.

\begin{figure*}[t]
\centering
\includegraphics[width=\textwidth]{fig2_composition.pdf}
\caption{Dataset composition, narrowing from what was assembled to what is
scored. (a) Evidence sources over the \num{74328} assembled variants; the
computational sources are near-saturating by construction while the clinical
ones are sparse. (b) The labelled subset on a log scale: the DMS pool outweighs
the clinical labels by a factor of \num{24}, and all of it comes from a single
\gene{MSH2} assay, which is why it cannot be used under leave-one-gene-out.
(c) The 683 clinical labels that every model result in this manuscript rests on,
split by gene, label source and class. \gene{MSH2}'s benign class is unlike the
other three: it carries \num{144} of the panel's \num{180} ProteinGym-clinical
benign calls against \gene{MLH1}'s 6, and those are the variants the error
analysis of Section~3.7 identifies as the source of every false positive.}
\label{fig:composition}
\end{figure*}

\subsection{Held-out-gene performance}

The grid completed all 16 cells without failure in 17.84\,h; all numbers below
are from the re-run described in Section~2.4, after the per-split seeding defect
was corrected. Holding the feature set at \texttt{esm+priors} and the PLLR mode
at \texttt{residual}, and averaging over three seeds and four held-out genes,
full fine-tuning did not improve ranking over the frozen backbone
(Table~\ref{tab:main}, Figure~\ref{fig:main}): AUROC $0.9283$ (SD $0.0111$) against $0.9335$
(SD $0.0100$), a difference of $-0.0052$, with $\Delta$AUPRC $-0.0039$. Both are
smaller than the frozen arm's own seed-to-seed spread. Excluding the unscoreable
\gene{PMS2} fold widens the AUROC difference to $-0.0152$ without changing its
sign. Fine-tuning does raise MCC, from $0.5495$ (SD $0.0988$) to $0.6553$
(SD $0.0241$), which is an operating-point effect rather than a ranking one and
is taken up in Section~3.4.

The feature branch moves performance more than freeze depth, and by a wider
margin than we previously reported. At seed 42, adding the 27 prior columns
raised AUROC from $0.9042$ to $0.9449$ on the frozen model ($+0.0407$) and from
$0.9270$ to $0.9409$ when fine-tuned ($+0.0139$). An earlier internal analysis
compared a fine-tuned model reading no priors against a frozen probe reading
them, observed a 6.5-point gap and attributed it to freeze depth; with the
factors separated, freeze depth accounts for approximately $0.005$--$0.015$
AUROC and the feature set for $0.014$--$0.041$. The two intervals do not
overlap, which is the substance of the correction: on this cohort the feature
set is the larger factor by roughly a factor of three, and the direction of the
freeze-depth effect is the opposite of what was originally claimed.

The most defensible reading is that at this training-set size the pretrained
representation is already close to what the head can exploit, so further
gradient updates largely fit the 380 training examples. That is a hypothesis
consistent with the observed divergence --- training loss fell roughly ninefold
while inner-validation AUROC peaked at epoch 2 and declined --- not something
these data establish.

\begin{table*}[t]
\caption{Main evaluation under leave-one-gene-out. Means over scoreable
held-out genes (\gene{MLH1}, \gene{MSH2}, \gene{MSH6}); three seeds for the two
ESM arms, one for the baseline. Thresholds selected on the inner validation
fold only; bold marks the best value in each column. \gene{PMS2} is excluded
throughout: 21 variants with 4 negatives fails the availability rule
(Section~2.4). Ranking (AUROC, AUPRC) is indistinguishable between the two ESM
arms while the operating point is not, which is why accuracy and AUROC are
reported together rather than either alone. The curated baseline, which reads
no language-model input, leads nine of the eleven columns.}
\label{tab:main}
\centering
\scriptsize
\setlength{\tabcolsep}{2.5pt}
%% Twelve columns at \small overran the text block by about 69pt and the last
%% three metrics were clipped off the page. This is as large as eleven metric
%% columns fit; do not enlarge either setting without re-reading the log for
%% "Overfull \hbox" against this table.
\begin{tabularx}{\textwidth}{@{}>{\raggedright\arraybackslash}X*{11}{c}@{}}
\toprule
Model & AUROC & AUPRC & Acc. & F1$_\text{path}$ & Bal. & Prec. & Rec. & Spec. & MCC & Brier & ECE \\
\midrule
ESM-2 650M frozen $+$ priors    & 0.9318 & 0.9376 & 0.6977 & 0.6966 & 0.7501 & 0.8411 & 0.7022 & 0.7980 & 0.5145 & 0.1919 & 0.2137 \\
ESM-2 650M full FT $+$ priors   & 0.9198 & 0.9306 & 0.8402 & 0.8442 & 0.8232 & 0.8515 & \textbf{0.8715} & 0.7749 & 0.6608 & 0.1501 & \textbf{0.1888} \\
Curated priors only, no ESM     & \textbf{0.9507} & \textbf{0.9442} & \textbf{0.8583} & \textbf{0.8467} & \textbf{0.8595} & \textbf{0.8858} & 0.8377 & \textbf{0.8813} & \textbf{0.7124} & \textbf{0.1374} & 0.1974 \\
\bottomrule
\end{tabularx}
\end{table*}

\paragraph{The baseline is the strongest model reported here.}
A curated-priors-only head, reading no ESM input, was run under the same
leave-one-gene-out protocol with \num{10000} bootstrap resamples and reached
AUROC $0.9651$ (95\% CI $0.9132$--$0.9957$) on \gene{MLH1}, $0.9078$
($0.8731$--$0.9386$) on \gene{MSH2}, $0.9791$ ($0.9520$--$0.9992$) on
\gene{MSH6} and $1.0000$ on \gene{PMS2} --- a mean of $0.9507$ over the three
scoreable genes, higher than any ESM-2 cell in the grid.

This comparison was previously withheld on the grounds that the baseline had
been scored on an earlier dataset build. That objection does not survive
inspection and is withdrawn. Re-running the baseline on the pinned build
reproduced every metric to the last recorded digit --- AUROC, AUPRC, MCC and all
bootstrap bounds on all four genes --- because the two builds expose the same 27
prior columns to this arm, including the four allele-frequency columns, over the
same clinical cohort. The gnomAD rebuild changed nothing this head reads, so the
arms were comparable throughout.

Two caveats attach to the row. It is a single deterministic run, while the two
ESM rows are three-seed means. Its threshold-dependent columns were computed
from the run's own per-variant predictions, selecting the threshold on the
inner-validation fold exactly as the grid cells do, and reproduce the run's
AUROC, MCC and per-gene thresholds to six decimal places. As with every other
arm, \gene{PMS2} is excluded from the reported mean; the $1.0000$ its results
file records over 21 variants is an artifact of cohort size, not a finding.

\begin{figure*}[t]
\centering
\includegraphics[width=\textwidth]{fig3_main.pdf}
\caption{Held-out-gene discrimination for the three headline arms, with 95\%
percentile-bootstrap confidence intervals (\num{10000} resamples). Filled marks
are seed 42; open marks are the seed-43 and seed-44 replicates of the two ESM
arms, shown to give the seed spread against which the between-arm differences
should be read. \gene{PMS2} is shaded: with 21 variants and 4 negatives it fails
the availability rule and is excluded from every mean.}
\label{fig:main}
\end{figure*}

\subsection{Calibration and threshold transfer}

Calibration is poor for every arm (ECE $0.214$ frozen, $0.189$ fine-tuned and
$0.197$ for the curated baseline; Brier $0.192$, $0.150$ and $0.137$), and no
post-hoc calibration was applied for the reason given in Section~2.4. The
baseline's calibration is no better than the language models': its ranking is
stronger, its probabilities are not.

Threshold transfer is the larger problem. Averaged over 16 cells and three
genes, the inner-validation MCC-optimal threshold yielded accuracy $0.7926$
against $0.7858$ for a fixed $0.5$ cut --- an advantage of $0.0068$ --- and the
fixed cut was better in 6 of 16 cells. In the worst case the selected threshold
produced accuracy $0.585$ against an achievable $0.912$, in a cell whose AUROC
is a healthy $0.938$. Fitting the best possible threshold on the held-out gene
itself, a deliberate leak reported only as a diagnostic ceiling, gives $0.917$
for the best cell and $0.883$ on average against $0.793$ achieved --- roughly
$0.09$ accuracy on average, and up to $0.33$ on individual cells, lost to
threshold placement rather than to ranking.

The mechanism is structural rather than a coding error: under leave-one-gene-out
the threshold is chosen on an inner fold drawn from the \emph{training} genes
and applied to a held-out gene whose score distribution differs, and the inner
fold holds roughly 95 variants, so an MCC-optimal cut estimated on it is noisy.
This argues for calibrating scores to be comparable across genes before a
threshold is transferred, rather than for a better threshold search.

\subsection{Per-gene robustness}

Mean AUROC by held-out gene across all 28 cells was $0.9201$ (SD $0.0299$) for
\gene{MLH1}, $0.8970$ ($0.0155$) for \gene{MSH2}, $0.9282$ ($0.0401$) for
\gene{MSH6} and $0.9291$ ($0.0504$) for \gene{PMS2}.

\gene{MSH2} is the lowest and the tightest, and Section~3.7 shows this is a
property of its labels rather than of the gene: it carries \num{144} of the
panel's \num{180} unreviewed ProteinGym-clinical benign calls, and every false
positive the headline arm makes is an \gene{MSH2} variant.

\gene{PMS2} results should not be reported. Its held-out set is 21 variants,
17 pathogenic and 4 benign, and it fails the availability rule in all 28 cells.
Its AUROC ranges from $0.7868$ to $0.9853$ --- a spread of $0.207$, the widest of
any gene and more than three times \gene{MSH2}'s --- because one misordered
variant moves the statistic by about $0.06$. The value of $1.0000$ recorded for
\gene{PMS2} in the priors-only baseline is an artifact of cohort size, not a
finding. \gene{PMS2} belongs in a coverage table, not a performance one; the one
place in this manuscript where including it would have changed a conclusion is
the allele-frequency increment of Section~3.6, and there we report the
three-gene answer and withdraw the claim.

Excluding \gene{PMS2} does not change the freeze-depth conclusion of
Section~3.3, so that finding is robust to its removal.

\begin{figure*}[t]
\centering
\includegraphics[width=\textwidth]{fig6_pergene.pdf}
\caption{Per-gene AUROC across all 28 cells (16 grid cells and 12
feature-family ablation cells, grey) with the three headline arms marked, and the held-out cohort
size beneath each gene. The spread within a gene is comparable to the spread
between arms, which is why no cell is ranked against another on a single gene.}
\label{fig:pergene}
\end{figure*}

\subsection{Ablations}

Table~\ref{tab:ablation} and Figure~\ref{fig:ablation} summarise the grid
together with four feature-family removals. Removing a feature group does not by itself remove every proxy for it:
AlphaMissense was trained on population and clinical data, so a ``without
gnomAD'' ablation is uninterpretable while AlphaMissense remains in the feature
set. The code enforces this, refusing an ablation that leaves a proxy for the
family being removed, which is why rows 5 and 7 are a joint bound and a
score-family removal rather than a clean allele-frequency arm. Rows 4--7 were
run at the grid's frozen \texttt{esm+priors} cell with every other flag held
identical, so that only the feature set moves.

\begin{table*}[t]
\caption{Ablation study. All cells share the dataset
(SHA-256 \texttt{78eb5d60\ldots}), split definition and evaluation code, and are
trained on ClinVar $+$ ProteinGym-clinical labels. Two AUROC columns are given
because the two conventions differ materially on \gene{PMS2}, whose 21-variant
fold is not scoreable: \textbf{4 genes} is the mean over all held-out folds,
\textbf{3 genes} the mean over \gene{MLH1}, \gene{MSH2} and \gene{MSH6} only.
Values are mean $\pm$ SD over seeds 42--44 where three seeds were run and
single-seed (42) otherwise; $\Delta$ is against row 3 frozen on the
three-gene convention. Rows 4--7 vary the feature set alone, at the frozen cell
of row 3. Row 9 was not run and is confounded by design.}
\label{tab:ablation}
\centering
\footnotesize
\setlength{\tabcolsep}{3pt}
%% The trailing Notes column was clipped off the page. Ablation and Notes share
%% what the six fixed columns leave, weighted so neither is starved.
\begin{tabularx}{\textwidth}{@{}c L{0.95} c c c c c L{1.05}@{}}
\toprule
\# & Ablation & Priors & Seeds & AUROC (4 genes) & AUROC (3 genes) & $\Delta$ & Notes \\
\midrule
1 & Curated tabular only     & 27, no ESM & 1 & 0.9630 & \textbf{0.9507} & $+0.0197$ & Best arm reported; Section 3.3 \\
2 & ESM-2 embedding only     & 0  & 1 & 0.9042 / 0.9270 & 0.9017 / 0.9125 & $-0.0293$ & frozen / full, seed 42 \\
3 & ESM-2 $+$ curated priors & 27 & 3 & $0.9335 \pm 0.0100$ & $0.9310 \pm 0.0055$ & --- & frozen; comparator for 4--7 \\
  &                          & 27 & 3 & $0.9283 \pm 0.0111$ & $0.9158 \pm 0.0092$ & --- & full fine-tune \\
4 & No structural features   & 25 & 3 & $0.9360 \pm 0.0105$ & $0.9277 \pm 0.0086$ & $-0.0032$ & pLDDT, disorder; within noise \\
5 & No gnomAD \emph{and} scores & 5 & 3 & $0.8804 \pm 0.0120$ & $0.8847 \pm 0.0161$ & $\mathbf{-0.0463}$ & Joint bound; proxy guard \\
6 & No UniProt/InterPro      & 24 & 3 & $0.9356 \pm 0.0051$ & $0.9337 \pm 0.0069$ & $+0.0028$ & Domains; within noise \\
7 & No external prior scores & 9  & 3 & $0.9075 \pm 0.0012$ & $0.8931 \pm 0.0072$ & $\mathbf{-0.0379}$ & 18 columns removed \\
8 & Frozen vs.\ fine-tuned   & 27 & 3 & $\Delta -0.0052$ & $\Delta -0.0152$ & --- & Section 3.3; primary result \\
9 & Label-source sensitivity & --- & --- & \todo{not run} & & & DMS pool single-gene; confounded \\
\bottomrule
\end{tabularx}
\end{table*}

\paragraph{Reading rows 4--7 against seed noise.}
Each ablation arm and its comparator are three-seed means, so the relevant scale
is the standard error of their difference, which ranges from $0.005$ to $0.010$
on the three-gene mean depending on the arm's own spread. Removing structural
features (row 4) moves the mean by $-0.003$ and removing domain features (row 6)
by $+0.003$ --- $0.55$ and $0.54$ standard errors respectively, the second in the
wrong direction to be a contribution at all. Neither family does detectable work
at this cohort size. Removing the external prior scores (row 7) costs $0.038$,
$7.2$ standard errors, and removing them together with allele frequency (row 5)
costs $0.046$, $4.7$.

\paragraph{Allele frequency's separate contribution does not survive on the
scoreable genes.} An earlier version of this table reported the increment of
row 7 over row 5 --- what allele frequency adds once the external scores are
already gone --- as $+0.0211$ against a standard error of $0.0052$, and read it
as a real if smaller effect. Recomputed after the seeding fix
(Section~2.4) the increment is $+0.0083$ against a standard error of $0.0102$,
which is indistinguishable from zero.

The reason is worth stating rather than burying, because it is a general hazard
of small held-out folds. On all four genes the same increment is $+0.0271$
($3.9$ standard errors) and would read as a clear effect. The entire difference
between the two conventions is \gene{PMS2} --- 21 variants, 4 of them negative,
the fold this manuscript declares unscoreable everywhere else. An effect that
exists only in the fold one has excluded on principle is not an effect one may
report, so we withdraw the claim. Rows 5 and 7 also differ by two feature
families, so even the four-gene figure would be the joint bound's remainder
rather than a clean allele-frequency estimate.

The prior block's contribution is therefore concentrated in 18 of its 27 columns
--- AlphaMissense, the zero-shot model scores and their rank transforms.
Structural annotation, domain annotation and, on the evidence available here,
allele frequency contribute nothing detectable once those 18 are present.

\paragraph{Why three seeds were necessary, and what seeding them properly changed.}
The comparator's three draws span $0.9253$ to $0.9364$ on the three-gene mean, a
range of $0.011$ (SD $0.0055$) --- still comparable to the row 4 and row 6
effects, so a single-seed table would have ranked those arms on a spread larger
than the differences it reported. Before the seeding fix the same three draws
spanned $0.027$ (SD $0.0142$). Roughly three-fifths of what this manuscript
previously attributed to seed variance was uncontrolled head initialisation:
the first fold of every run drew its classifier weights from process entropy
rather than from the seed (Section~2.4). Correcting it did not change which
ablations matter, but it more than halved the noise floor they are read against
--- and, as above, it removed one claim that had been sitting just above the old
floor. This is also why no cell is ranked against another on a single gene
(Figure~\ref{fig:pergene}).

\begin{figure*}[t]
\centering
\includegraphics[width=\textwidth]{fig5_ablation.pdf}
\caption{Every arm, seed-averaged, ranked by mean AUROC over the scoreable genes
against the curated-features-only baseline (horizontal line). Bars are $\pm 1$
SD over seeds where three were run; single-seed arms are plotted as points and
labelled as such, so a one-draw cell is never read as a measured mean. The two
ablations that remove structural or domain features sit level with the
comparator; the two that remove external prior scores sit at the bottom, below
every cell that keeps them. No configuration reaches the baseline.}
\label{fig:ablation}
\end{figure*}

Two further axes --- PLLR mode and fusion --- were varied at a single seed each.
We do not interpret them: the three-seed axis shows an MCC standard deviation up
to $0.131$, large enough that single-seed differences should not be read as
effects.

\subsection{Error analysis}

Errors were taken per seed, at the threshold each fold selected on its own inner
validation, and intersected across the three seeds of an arm. A variant that one
seed calls wrong and two call right is a draw of the initialisation rather than a
property of the variant, and \num{193} of the \num{683} held-out variants are of
that kind for the frozen arm --- a per-variant view of the seed spread quantified
in Section~3.6. Only the \num{85} variants all three seeds get wrong are counted
below.

\paragraph{The errors are concentrated in one gene and one label source.}
The frozen arm makes \num{85} consensus errors on \num{683} held-out variants
($12.4\%$): \num{56} false positives and \num{29} false negatives. Every one of
the \num{56} false positives is a \gene{MSH2} variant. \gene{MLH1} and
\gene{MSH6} produce none at all; their errors are exclusively false negatives.
The full fine-tune behaves the same way with more errors --- \num{97} total,
\num{73} false positives, \num{68} of them \gene{MSH2}.

Testing every covariate the master table carries, and separating those the model
reads from those it does not, two stand out and neither is a model input:
\texttt{label\_source} and \texttt{evidence\_tier} (Benjamini--Hochberg
$p = 1.0 \times 10^{-4}$ for the frozen arm, $1 \times 10^{-6}$ for the fine-tuned
one). This separation matters: an association with a feature the model reads
would restate the model's own weighting rather than tell us anything about it.
\num{45} of the \num{56} false positives are ProteinGym-clinical benign calls,
and \num{44} carry \texttt{evidence\_tier == unreviewed}. The model assigns these
probabilities above $0.999$.

\paragraph{This explains the \gene{MSH2} operating-point anomaly.}
\gene{MSH2} has been the outlier in every table in this manuscript --- respectable
ranking (AUROC $0.897$ averaged over the grid) alongside the weakest
threshold-dependent performance of any scoreable gene. Figure~\ref{fig:composition}(c)
shows why. \gene{MSH2} contributes \num{144} of the panel's \num{180}
ProteinGym-clinical benign labels; \gene{MLH1} contributes \num{6}. Its benign
class is therefore composed largely of calls that carry no review status, and
those are precisely the variants the model scores as pathogenic with high
confidence. The gene is not harder in any biological sense we can demonstrate;
its negative class is annotated differently from the other three.

\paragraph{What this does not license.} The association is with label provenance,
not with biology. These \num{45} variants may be mislabelled, or they may be
genuinely benign variants that every feature in the block scores as damaging;
distinguishing those requires returning to the primary submissions, which is
outside the scope of this work. We therefore make no claim about individual
variants, and we do not report a ``high-confidence false positive'' list as a
candidate reannotation set. What the analysis does establish is narrower and
more useful: on this panel, held-out-gene performance is partly a measurement of
label provenance, and a cohort assembled with uniform review standards would be
required to separate the two. Per-variant exports are in
\texttt{error\_analysis\_<arm>\_errors.csv} and the covariate tests, with the
model-input flag on each, in \texttt{error\_analysis\_<arm>\_covariates.csv}.

% =====================================================================
\section{Discussion}
% =====================================================================

The substantive result is methodological. When freeze depth and feature set are
varied independently, freeze depth matters little for ranking on this cohort
while the curated prior block matters more. This corrects an earlier reading of
our own data in which a 6.5-point gap was attributed to freeze depth by
comparing a fine-tuned model that read no priors against a frozen probe that
did. We report the correction because the confound is easy to reproduce: any
comparison between a fine-tuned language model and a feature-based probe must
hold the feature set constant, or it measures both at once.

The stronger reading is that on this cohort the language model did not earn its
place at all. A head reading 27 curated columns and no ESM input ranks better
than every configuration that adds a 650M-parameter backbone, on nine of eleven
reported metrics. We do not present this as a general result about protein
language models --- 683 clinical labels across four genes cannot support one ---
but it does set the bar that a language-model contribution has to clear here,
and none of the 16 configurations cleared it. The ablations then locate most of
the curated signal in features that are themselves the outputs of other
predictors, which is a weaker claim about biology than the raw figure suggests.

Second, the threshold protocol accounts for much of the difference between the
arms as a user would experience it. The two ESM models are indistinguishable in
AUROC and AUPRC yet differ by $0.143$ in accuracy and $0.169$ in recall. Under
leave-one-gene-out this is expected, and the diagnostic ceiling suggests roughly
$0.09$ accuracy is recoverable by calibration. We could not test this, because
the pipeline discarded the predictions a calibrator would need.

Held-out-gene evaluation is why these numbers are as low as they are and why we
consider them meaningful. We did not run a random-split comparison, so we cannot
quantify the inflation on this data; that is an explicit gap rather than a
rhetorical claim.

\paragraph{Limitations.}
The supervised cohort is 683 clinically labelled variants across four genes, one
of which cannot be scored. Confidence intervals on individual genes are
correspondingly wide --- the priors-only baseline's \gene{MLH1} interval spans
$0.913$ to $0.996$. The DMS pool, though large, is one assay on one gene and is
therefore not a general supervision source under a gene-held-out protocol.
Calibration is poor for every model reported. ClinVar labels carry ascertainment
bias: variants are submitted because someone had reason to test them, so the
labelled set is not a random sample and the benign class is enriched for common
alleles. The panel is four genes of one pathway in one disease.

\paragraph{Residual circularity.}
AlphaMissense is in the feature set and was trained on ClinVar-derived data, so
part of the measured agreement with ClinVar labels may reflect agreement with an
upstream model. The code quarantines frequency-derived priors when frequency
mints labels and drops gene-constant priors under leave-one-gene-out, but it
does not decorrelate AlphaMissense from ClinVar. Ablation 7 is the direct test,
and its result supports the concern rather than dispelling it: removing the 18
external score columns costs more AUROC than removing any other feature family,
about four standard errors, while removing structure and domain annotation costs
nothing measurable. The curated block's contribution is therefore concentrated
in features that encode other predictors' opinions. We cannot separate how much
of that is independent evidence from how much is agreement with models trained
on overlapping data; doing so needs a label set disjoint from AlphaMissense's
training data, which this cohort does not provide.

\paragraph{This is not a clinical tool.}
It is a research benchmark trained on 683 labelled variants from four genes,
evaluated on cohorts as small as 21, with expected calibration error near
$0.19$. No threshold reported here has been validated against clinical outcomes.
Before clinical consideration the minimum steps are a substantially larger
multi-centre labelled cohort; prospective evaluation on variants absent from all
training sources; calibration against observed pathogenicity rates rather than
held-out labels; per-population analysis, which the present data cannot support;
and evaluation of the specific decision the tool would inform.

% =====================================================================
\section{Conclusion}
% =====================================================================

Separating backbone freeze depth from feature set removes an effect we had
previously attributed to freeze depth: full fine-tuning of a 650M-parameter
protein language model does not improve ranking over a frozen backbone on this
cohort, while the curated prior block does. More of the practical difference
between models arises from threshold transfer across held-out genes than from
the models themselves, and that is the most tractable target for improvement.
The absolute figures rest on 683 clinical labels across four genes with one gene
unscoreable, and should be read as a benchmark for method development rather
than as evidence of clinical utility. The comparison that matters most ---
protein language model against curated features alone --- goes against the
language model: a curated head with no ESM input ranks above every cell in the
grid, and feature-family ablations place most of that curated signal in columns
inherited from other predictors rather than in structural or domain evidence.
The most tractable target for improvement is not model capacity but the
threshold and calibration protocol that transfers across held-out genes.

% =====================================================================
\section*{Acknowledgements}
\todo{Acknowledgements}

\section*{Funding}
\todo{Funding statement --- do not invent funders or grant numbers}

\section*{Conflict of Interest}
\todo{Declaration}

\section*{Data availability}
Build manifests with per-source URLs and SHA-256 checksums, all 28 cells'
per-variant predictions, run summaries and hardware/environment captures are in
the repository. The table underlying every model result reported here is
identified by SHA-256
\texttt{78eb5d60860cc08e\allowbreak adb6faa0c0d7fbd2\allowbreak
2adbe6f277bb95e9\allowbreak 7a3a680264b4430d}. Grid
cells were run without checkpoint saving, so backbone weights are not
distributable. AlphaMissense is CC BY-NC-SA 4.0; \todo{confirm redistribution
terms for any derived table embedding AlphaMissense values}.

\bibliographystyle{unsrt}
\bibliography{references}

\end{document}
